\section{Arithmetic and the class-two comparison}\label{sec:abelian}

This section supplies the arithmetic input and the cup comparison
for Corollary~\ref{thm:class-two-components}.
Write $\Delta=\Gal(F_1/\Q)$, choose the Teichm\"uller character
$\omega$, and read character exponents modulo $p-1$.
For a compact $\Z_p[\Delta]$-module $V$, put
\[
 V^{[j]}=\epsilon_jV,\qquad
 \epsilon_j=\frac1{p-1}\sum_{\delta\in\Delta}\omega(\delta)^{-j}\delta.
\]
Let $\Gamma=\Gal(F_\infty/F_1)$, choose $\gamma$ with
$u=\chi_p(\gamma)$ generating $1+p\Z_p$, and put
$\Lambda=\Z_p[[\Gamma]]$, $R=\Lambda[1/p]$.
The abelianization $M=H^{\rm ab}$ has its conjugation action of
$\Delta\times\Gamma$, independently of choices of lifts. Put
\[
 U=M^-,\quad T=M^+=\bigoplus_{e\ {\rm even}}T_e,\quad
 T_e=M^{[e]},\quad D=U_{\Lambda\text{-tor}}.
\]

Choose a lift of $\Delta$ by profinite Schur--Zassenhaus.
Coprime fixed-point lifting gives an element of
$\operatorname{Gal}(\Omega_p/F_1)^\Delta$ above a topological
generator of $\Gamma$. Its procyclic closure maps isomorphically
to $\Gamma$, giving commuting lifts and a splitting
$\operatorname{Gal}(\Omega_p/F_1)=H\rtimes\Gamma$.
The lifting assertion follows at each finite pro-$p$ quotient
from Schur--Zassenhaus applied to the preimage of the fixed
coset; compactness gives compatible lifts. These choices are
used for the cochain calculation. The full invariance calculation
in Section~\ref{sec:class-two-components} removes them from the
resulting quotient.

\subsection{Norms and the odd torsion}

The following computation removes the use of finiteness of the whole
plus module in the component theorem. It is a fixed-coefficient
form of the unit-radical phenomenon in
\cite[Theorem~3, p.~3]{McCallum}.
\begin{lemma}\label{lem:norms}
Put $P_r=H\rtimes\Gamma^{p^r}$ and $A_n=\Z/p^n(1)$.
Let $D_0\subseteq H^{\rm ab}$ be a $\Gamma$-stable finitely
generated $\Z_p$-submodule. A sequence
$f_r\in H^1(P_r,A_n)$ compatible under corestriction restricts
to zero on $D_0$ at every level where it is defined.
\end{lemma}
\begin{proof}
Restriction to $H$ factors through $H^{\rm ab}$, because $H$
acts trivially on the coefficients. Write
$\phi_r=(\operatorname{res}_Hf_r)|_{D_0}$ and
$V_n=\Hom_{\Z_p}(D_0,A_n)$. On this finite module $\gamma$ acts by
\[
                  (B\phi)(d)=u\phi(\gamma^{-1}d).
\]
The corestriction formula with coset representatives
$\gamma^{jp^r}$, $0\leq j<p$, gives
\[
 \phi_r=N_r(B)\phi_{r+1},\qquad
 N_r(B)=\sum_{j=0}^{p-1}B^{jp^r}.
\]
Continuity gives an $s$ with $B^{p^s}=1$ on $V_n$.
Hence $N_r(B)=p$ for $r\geq s$. Iterating $n$ times gives
$\phi_r=p^n\phi_{r+n}=0$ at all sufficiently high levels.
The transition formulas prove the assertion at earlier levels.
\end{proof}

For $p\nmid a$ the identity
\[
 N_{F_{r+1}/F_r}(1-\zeta_{p^{r+1}}^a)=1-\zeta_{p^r}^a
\]
provides such a sequence of Kummer classes at every fixed coefficient
exponent $n$. For $a=p^vb$, $p\nmid b$ and $v<r$, one has
\[
 1-\zeta_{p^r}^{p^vb}
   =\prod_{j=0}^{p^v-1}(1-\zeta_{p^r}^{b+jp^{r-v}}).
\]
All exponents on the right are prime to $p$. Products and the
$c$-even projector preserve corestriction compatibility.
Every cyclotomic unit class used below consequently annihilates $D$.

The group $H$ is free pro-$p$ by Ferrero--Washington and
Iwasawa \cite[Theorem~2, pp.~273--274]{Iwasawa}.
Its abelianization $M$ is finitely generated over $\Lambda$:
remove the $\Gamma$-coordinates from a finite generating set of
$\Gal(\Omega_p/F_1)$ and take their abelian classes in $H$.
Their $\Gamma$-conjugates generate $M$.
In particular $M$ is $\Z_p$-torsionfree.
The module $D$ is finite free over $\Z_p$. To see this, choose
an annihilator $p^af$ of its finitely generated $\Lambda$-torsion
module, with $p\nmid f$. Injectivity of multiplication by $p$
gives $fD=0$. Weierstrass preparation makes $\Lambda/(f)$ finite
over $\Z_p$. Moreover $D=\ker(f:U\to U)$, so it is closed.
It is saturated: $px\in D$ gives $fpx=0$, then $fx=0$, hence
$x\in D$.

Put $W=U/D$. Local and global unit duality gives $\Lambda$-rank
one in every odd character of $W$
\cite[Proposition~2.2.6 and its proof, pp.~3--4]{WWEOrd}.
The global minus units contribute roots of unity, and the
unramified quotient has rank zero. The ring $R$ is a PID:
$\Lambda$ is a two-dimensional regular local UFD, localization
at $p$ removes its height-two prime, and a one-dimensional
noetherian UFD is a PID. Therefore
\begin{equation}
                  W[1/p]=\bigoplus_{a\ {\rm odd}}Re_a .
                  \label{eq:Wfree}
\end{equation}
Finite generation clears the denominators of this isomorphism
and its inverse. Completed tensor powers thus have bounded
denominators. The exterior square has alternating coordinates
in $\Z_p[[T_1,T_2]][1/p]$.

\subsection{The two arithmetic characters}

The arithmetic inputs occur in two explicitly related components.
The local/global unit sequence in
\cite[Corollary~2.1.3, p.~8]{WWEPseudo} has the form
\begin{equation}
 0\longrightarrow\mathcal E_e/\mathcal C_e
 \longrightarrow\Lambda/(f_e)\longrightarrow T_e
 \longrightarrow X^{[e]}\longrightarrow0 ,
 \label{eq:unitseq}
\end{equation}
with equal characteristic ideals for $\mathcal E_e/\mathcal C_e$
and $X^{[e]}$. Here $\mathcal E_e,\mathcal C_e$ are the norm limits
of global and circular units in character $e$.
Finiteness of $X^{[e]}$ kills both end terms after inverting $p$;
hence $T_e[1/p]$ is cyclic.
The even $p$-ramified module $T$ is finite free over $\Z_p$ by
Kummer duality, Ferrero--Washington, and the torsionfreeness of $M$
\cite[Proposition~3.4.7 and Corollary~3.4.8, p.~92]{SharifiIwa}.
Let $q_e(z)$ be its characteristic polynomial on $T_e$ for
$\gamma$; cyclicity makes it the minimal polynomial as well.

The cup dual has action
\begin{equation}
 (g\lambda)(t)=\chi_p(g)^2\lambda(g^{-1}t).
 \label{eq:twodual}
\end{equation}
Thus $T_e^*(2)$ has character $\omega^{2-e}$; an eigenvalue
$\alpha$ becomes $u^2/\alpha$. Kummer adjunction identifies
$T_e[1/p]$ with the corresponding dual and twist of
$X^{[1-e]}[1/p]$, using the analogous formula with one twist.
The symbol theorem therefore uses
$\theta=\omega^{2-e}$ and $\chi=\omega^{1-e}$.
The excluded characters $\theta=1,\omega^2$ correspond to
$T_2,T_0$, which are zero by $X^{[-1]}=X^{[1]}=0$
\cite[proof of Lemma~5.2.7, p.~13]{WWEOrd}.

\subsection{The one-component cup input from Ohta's isomorphism}
\label{subsec:ohta-component-gate}

We specify the published isomorphism used here, including its recently
completed density input.  Ohta's Part~II, Theorem~(6.2.2), p.~409, in \cite{Ohta2020} proves the isomorphism given
by the upper-right coordinate of the Eisenstein Galois representation.
Its use of CM-point density in Part~I, p.~381, follows from Ohta's \cite[Theorem~(2.3.4), author pp.~15--16]{OhtaDensity}.  This match can be checked directly: the shifts
$\delta_{\mathfrak q}=\{\mathfrak q\}_\Gamma$ chosen on p.~378 of the
2020 paper lie in $\ker(\mathrm{Cl}_\infty\to\mathrm{Cl}_{n_1})$ and
are distinct modulo $\mathrm{Cl}^{\rm alg}$ by its
Proposition~(7.2.4)(iv).  They, together with the unbounded conductor
subsequence $n(\epsilon)$ and the classes in
$\ker(\mathrm{Cl}_{n(\epsilon)}\to\mathrm{Cl}_{n_1})$, are exactly the
data of the 2026 theorem with initial level $n_1$.  Forgetting the
finite level structures maps the former CM tuples onto the latter.
Finiteness and irreducibility of the modular cover then transfer the
density back to the source used in 2020.  Thus the required mod-$p$
nonvanishing input is a theorem.

Here is the action dictionary.  For an even nonexceptional $e$, set
\[
 \chi=\omega^{1-e},\qquad
 \theta_{\mathrm{WWE}}=\omega^{2-e},\qquad
 \theta_{\mathrm O}=\omega^{e-2}.
\]
We define the modular target separately from the arithmetic module
$M=H^{\mathrm{ab}}$.  In the modular-curve and dual-Hecke convention of
\cite[\S3.1]{WWEPseudo}, put $\theta=\theta_{\mathrm{WWE}}$ and
\[
 H'_{\mathrm{mod},\theta}
 =\varprojlim_r H^1_{\mathrm{et}}
   (X_1(p^r)_{\overline{\mathbb Q}},\mathbb Z_p)^{\mathrm{ord}}_\theta,
 \qquad
 H_{\mathrm{mod},\theta}
 =H'_{\mathrm{mod},\theta}\otimes_{\mathfrak h'_\theta}\mathfrak h_\theta.
\]
Here $X_1(p^r)$ is the compactified modular curve classifying
elliptic curves with a point of exact order $p^r$, and the ordinary
summand is the summand on which the dual $U_p$ operator is invertible.
The transitions are corestrictions, the subscript is the diamond
eigenspace, and $\mathfrak h_\theta$ is the localization of the algebra
$\mathfrak h'_\theta$ of dual Hecke operators at its Eisenstein maximal
ideal.  Let $I_\theta\subset\mathfrak h_\theta$ be the cuspidal Eisenstein
ideal, the image of the annihilator of the ordinary Eisenstein series
in the full Hecke algebra.  Set
\[
 M_{\mathrm{mod}}
 =H^-_{\mathrm{mod},\theta}/I_\theta H^-_{\mathrm{mod},\theta},
\]
where the minus sign denotes complex conjugation.  The coefficient
$\Lambda$ acts by $\gamma\mapsto\langle\kappa(\gamma)\rangle$.

Ohta's class character is
$(\theta_{\mathrm O}\omega)^{-1}=\chi$.  His tame level is $1$,
and the exponent $e-2$ is even, so that $-1$ belongs to its kernel,
as required in Theorem~(6.2.2).  The exceptional indices $0$ and $2$
have $T_0=T_2=0$ and need no cup-symbol argument.  Ohta's $X^-$ is a
splitting of the ordinary quotient $X/X^+$, with $X^+=X^{I_p}$
in his (6.1.14).  The comparison with the above minus space uses the
natural isomorphism $H^-_{\mathrm{mod},\theta}\to H_{\mathrm{quo},\theta}$
of \cite[Theorem~3.4.1 and Corollary~3.4.2]{WWEPseudo}, followed by the
modular-curve and diamond-action transport to that ordinary quotient.
In Ohta's (6.1.26), an element $\gamma$ with $\kappa(\gamma)=u$ acts on
$(X^-/I^*X^-)^\dagger$ as $u^{-1}\iota(u)^{-1}$.  After the Tate twist $(1)$,
this becomes the inverse diamond action $\iota(u)^{-1}$; its character
is $\theta_{\mathrm{WWE}}$.  The extension-coordinate map is
\[
 \sigma\longmapsto \sigma e^+-e^+\pmod{I^*X}.
\]
Sharifi's \cite[pp.~3 and 13--14]{SharifiProgram} explicitly uses the inverse-Galois convention and records the
Ohta isomorphism for this $\Upsilon$.  Equivalently, this is the
upper-right-coordinate definition in \cite[\S8.1]{WWEPseudo}:
evaluate the extension coordinate in
$\operatorname{Hom}_{\mathfrak h_\theta/I_\theta}
 (H^+_{\mathrm{mod},\theta}/I_\theta H^+_{\mathrm{mod},\theta},
 M_{\mathrm{mod}})$ on the distinguished element corresponding to
$1$ under $\mathfrak h_\theta/I_\theta\simeq
H^+_{\mathrm{mod},\theta}/I_\theta H^+_{\mathrm{mod},\theta}$.
Its restriction to the cyclotomic kernel is a homomorphism, and
ordinarity makes it factor through $X_\chi(1)$.
We therefore have the isomorphism
\[
 \Upsilon: V=X^{[1-e]}(1)\xrightarrow{\ \sim\ }M_{\mathrm{mod}}.
\]

Suppose now that $X^{[e]}$ is finite, so that the previously established
adjoint calculation makes $V[1/p]$ cyclic.  At a characteristic prime
$f$, let $R$ be the corresponding discrete valuation ring and write
$V_R\simeq R/(f^r)$.  Transport this generator to
$(M_{\mathrm{mod}})_R$ by $\Upsilon$.
The endomorphism $\Upsilon\varpi$ is multiplication by an element $a$,
and the Fukaya--Kato identity in
\cite[Theorem~5.21]{SharifiFK} gives
\[
 (a-1)\xi'=0\quad\hbox{in }R/(f^r).
\]
Writing $\xi=f^r g$ with $g\in R^\times$, separability in
characteristic zero gives $v_f(\xi')=r-1$.  Hence $a-1\in(f)$,
so $a$ is a unit.  Thus $\varpi_R$ is an isomorphism at every
characteristic prime, with no restriction on the multiplicity $r$.
It follows that the rational cup-symbol map is surjective on all of
$V[1/p]$. Both integral modules are finitely generated over $\Lambda$
with zero $\mu$-invariant, so the cokernel of the integral map
is a finite $\mathbb Z_p$-module. One power of $p$ kills it.
Reduction along the norm-compatible finite-level symbol maps
therefore gives the uniform integral bound used below.

Thus $X^{[e]}$ finite supplies cyclicity and the full rational
cup image in the twice-twisted companion.  The arithmetic character
transport is the displayed map $\theta_{\mathrm O}\mapsto
\theta_{\mathrm{WWE}}$ together with the Tate twist.

The cyclicity step is used in the direction displayed in the
Fukaya--Kato identity. For example, in a discrete valuation ring
with uniformizer $t$ and residue characteristic zero, take
$V=R/(t^2)$, $M_{\mathrm{mod}}=(R/t)^2$ and both maps zero. Then
$t(\Upsilon\varpi-1_{M_{\mathrm{mod}}})=0$, while $tV\ne0$.
In the argument above, Ohta's isomorphism identifies $M_{\mathrm{mod}}$ with the
cyclic module $V$ and turns the identity into multiplication by
$a$; that is the step which excludes this example.


For clarity we pass to group cohomology explicitly.
Put $S_n=\operatorname{Spec}\mathcal O_{F_{n+1}}[1/p]$.
The maximal pro-$p$ quotient of $\pi_1(S_n)$ is $P_n$:
normal closures of finite $p$-extensions of cyclotomic layers
remain $p$-extensions over $F_1$.
Inflation is an isomorphism on $H^1$ and an injection on $H^2$.
Indeed, an invariant homomorphism from the kernel of the
maximal pro-$p$ quotient to a finite $p$-module would give
a larger pro-$p$ quotient. The five-term sequence applies.
The universal finite \'etale tower has zero $H^1$ with finite
local coefficients, since every torsor becomes trivial in
that tower. Cartan--Leray consequently gives the same comparison
with \'etale cohomology \cite[Theorem~14.9 and
Proposition~14.13, pp.~96--98]{MilneEtale}.
These maps respect cups. The cup-symbol map just proved has
finite index in its integral target. Fix $p^{b_0}$ annihilating
that finite cokernel. The unique prime above $p$ is totally
ramified in every relative cyclotomic layer, so its intersection
with the lower Hilbert $p$-class field is trivial. The class norms
are therefore surjective. Consequently the projections
$V=X^{[1-e]}(1)\to H^2_{\mathrm{et}}(S_n,\Z/p^n(2))^{[2-e]}$
are onto, and $p^{b_0}$ times every such target lies in the
projection of the integral cup-symbol image. The comparison of
\cite[\S\S7.1.9--7.1.10, p.~86]{FK} then puts
$p^{b_0}$ times the asserted cohomology component in the image
of the cyclotomic-unit cups, uniformly at every level.

Reduction from coefficients $p^{n+1}$ to $p^n$ is surjective
by the class-group description in \cite[\S7.1.9]{FK}.
There is no Brauer contribution: the field is totally imaginary
and only one finite prime is omitted.
For a $p$-unit $v$, the quotient $v/c(v)$ is a global unit
with absolute value one at every embedding, hence a root of
unity. The odd unit classes are cyclic, and their alternating
self-cup is zero because $p$ is odd. Thus, for every $n\geq1$,
the subgroup $p^{b_0}H^2(P_n,\Z/p^n(2))^{[2-e]}$
is generated by cups of $c$-even classes annihilating $D$.
The exponent $b_0$ depends on $e$ and $p$, and is independent
of the level $n$.

\subsection{The projected class-two extension}

We give the comparison explicitly, because removing $D$ changes
the source of the cup calculation.
For a pro-$p$ group, brackets below mean closed commutator
subgroups. Since $p$ is odd, the quotient $H/[H,[H,H]]$
has its integral class-two Baker--Campbell--Hausdorff Lie
structure. Its central commutator module is
$\widehat{\bigwedge}^{\,2}_{\Z_p}M$. This follows first for a
finite free generating set by collecting commutators, and then
for $H$ by the inverse limit over finite subsets of a convergent
free generating sequence. Kill the central brackets involving
$T$. The remaining central module is
$\widehat{\bigwedge}^{\,2}U$, on which $c$ acts as $1$.
The minus part of this Lie module consequently maps
isomorphically to $U$; this gives a canonical odd section.
Kill its image of $D$, together with all brackets involving
$D$. The resulting Lie module has abelianization $W\oplus T$,
central module
\[
                   N=\widehat{\bigwedge}^{\,2}_{\Z_p}W,
\]
and plus part in the exact sequence
\begin{equation}
             0\longrightarrow N\longrightarrow \mathcal E
              \longrightarrow T\longrightarrow0 .
             \label{eq:projectedextension}
\end{equation}
All these quotients are stable under $\Gamma$ and $\Delta$.
The closedness assertions follow from compactness. The quotient
$W$ is $\Z_p$-torsionfree, so its completed exterior square
is torsionfree, as is seen in product $\Z_p$-bases.

Let $E$ denote the even exponents $e$ for which $X^{[e]}$
is finite. For $e\in E$ with $T_e\ne0$, choose a cyclic vector
$v_e\in T_e[1/p]$ and a lift $y_e\in\mathcal E^{[e]}[1/p]$.
Set
\begin{equation}
                C_e=q_e(\gamma)y_e\in N[1/p].
                \label{eq:projectedcompanion}
\end{equation}
The action of $\gamma$ on two odd coordinates is multiplication
by $z=(1+T_1)(1+T_2)$. Thus $C_e$ has alternating coefficient
series $C_{e;a,b}(T_1,T_2)$, with $a,b$ odd and
$a+b=e$ modulo $p-1$. Each series has a common bounded
denominator.

\begin{lemma}\label{lem:projectedpole}
For each root $\alpha$ of $q_e$, at least one matching coefficient
of $C_e$ has nonzero restriction to
$(1+T_1)(1+T_2)=\alpha$.
\end{lemma}
\begin{proof}
Let $L/\Q_p$ split $q_e$, with ring of integers $\mathcal O$.
Every eigenvalue satisfies $\alpha-1\in\mathfrak m_{\mathcal O}$,
since the action of $\Gamma$ on the finite lattice $T_e$ is
continuous. Suppose all the restrictions vanish. The coordinates
\[
 S=(1+T_1)(1+T_2)-\alpha,\qquad
 T_2=\frac{\alpha+S}{1+T_1}-1
\]
give an isomorphism of complete local power-series rings over
$\mathcal O$. Vanishing at $S=0$ is divisibility by $S$.
Division preserves bounded denominators and alternation.
It follows that $C_e=(\gamma-\alpha)A$ for some $A\in N_L$.

Write $r(z)=q_e(z)/(z-\alpha)$. In the cyclic module $T_{e,L}$,
the vector $w=r(\gamma)v_e$ is nonzero and spans an
$\alpha$-eigenline $J$. This remains true for a repeated root:
$r$ has degree one less than the annihilating polynomial.
The vector $r(\gamma)y_e-A$ lifts $w$ equivariantly.
Thus \eqref{eq:projectedextension} splits over $J$.

The lattice $J\cap T_{\mathcal O}$ is a direct summand of
$T_{\mathcal O}$: its quotient is torsionfree over the DVR
$\mathcal O$. Extend this equivariant lift to an
$\mathcal O$-linear section
$s:T_{\mathcal O}\to\mathcal E_{\mathcal O}[1/p]$.
Some fixed $p^b$ makes its image integral. For $h\in H$ write
its image in the class-two Lie module as
\[
        x(h)+s(t(h))+z(h),\quad
 x(h)\in W_{\mathcal O},\ t(h)\in T_{\mathcal O},\
 z(h)\in N_{\mathcal O}[1/p].
\]
The same $p^b$ makes every $z(h)$ integral. The section defect
\[
        \beta_\delta(t)=\delta s(t)-s(\delta t)
\]
vanishes on $J$ for $\delta\in\Gamma$.

Fix $n$ and take two of the even cyclotomic unit classes $f,g$
whose cups span the indicated finite-index subgroup. Their invariant cocycle
representatives are obtained by averaging with $(1+c)/2$.
They vanish on the $\Gamma^{p^n}$ section, since $c$ fixes
that section and negates $\mathcal O/p^n(1)$. On $H$ they
vanish on $T$ by parity and on $D$ by Lemma~\ref{lem:norms};
therefore they factor through $W$.
Let $\psi=f\wedge g:N_{\mathcal O}\to\mathcal O/p^n(2)$.
Define the continuous cochain
\[
 k(h)=-\frac{p^b}{2}f(h)g(h)-\psi(p^b z(h)),\qquad
 k(h\delta)=k(h).
\]
The class-two product adds $[x(h),x(h')]/2$ to its central
coordinate. Using $dk(h,h')=k(h')-k(hh')+k(h)$ on $H$
gives
\[
                       dk(h,h')=p^bf(h)g(h').
\]
Indeed the quadratic term contributes
$p^b(f(h)g(h')+f(h')g(h))/2$, and the central term contributes
$p^b(f(h)g(h')-f(h')g(h))/2$.
For $\delta,\delta'\in\Gamma^{p^n}$ the remaining cocycle is
\begin{equation}
 (p^bf\smile g-dk)(h\delta,h'\delta')
       =-\psi\bigl(p^b\beta_\delta(t(h'))\bigr).
 \label{eq:projectedcochain}
\end{equation}
To verify this formula, use equivariance of $f,g$ on $H$,
their vanishing on the section, and
$z(\delta h'\delta^{-1})=\delta z(h')+
\beta_\delta(t(h'))$. This also verifies the coefficient
twist $\chi_p(\delta)^2$ in the coboundary.
The edge class in \eqref{eq:projectedcochain} restricts to zero
on $J$.

We spell out why this contradicts the uniform finite-index assertion
without making compatible choices of cup expressions.
Both $H$ and $\Gamma^{p^n}$ have $p$-cohomological dimension
one. Hochschild--Serre gives
\[
 H^2(P_n,\mathcal O/p^n(2))^+
 \simeq H^1(\Gamma^{p^n},
             \Hom_{\Z_p}(T,\mathcal O/p^n)(2)).
\]
Put $V=\Hom_{\Z_p}(T,\mathcal O)(2)$ and let $B$ be its
$\gamma$-operator. Evaluation on the procyclic generator
identifies this group with
\begin{equation}
             V/(p^nV+(B^{p^n}-1)V).
             \label{eq:coinvariants}
\end{equation}
Corestriction sends a generator value to that same value
modulo $(B^{p^n}-1)V$: this follows by evaluating the
corestriction cocycle on $\gamma^{p^n}$, where exactly one
coset representative crosses a fundamental domain.
Thus these transitions are quotient maps.
Continuity makes their defining submodules cofinal with the
$p$-adic neighbourhoods of zero, so their inverse limit is $V$.

Restriction to $J\cap T_{\mathcal O}$ is a surjection on
integral dual lattices. Its target has character $2-e$ and
$\gamma$-eigenvalue $u^2/\alpha$ by \eqref{eq:twodual}.
It induces surjections on the finite quotients
\eqref{eq:coinvariants}. By the uniform finite-index assertion and
\eqref{eq:projectedcochain}, the image of the fixed submodule
$p^{b+b_0}V^{[2-e]}$ is zero in every such quotient of the dual
lattice of $J$. Their separated inverse limit is that
nonzero free $\mathcal O$-module. A surjection onto it cannot
annihilate $p^{b+b_0}V^{[2-e]}$. This contradiction proves the lemma.
\end{proof}

\subsection{The rational outer image and its odd coefficients}

We construct the rational target and prove that its linear
coefficients detect the quotient $W$.
Hain--Matsumoto's Theorem~8.4 and the surrounding construction,
together with Brown's Theorem~1.1, give the weight-completed
free Lie algebra on $s_m$ for odd $m\geq3$ as the Lie algebra
of the rational outer image \cite{HM,Brown}.
The image of $H$ is dense in every finite weight quotient.
The reason is that the cyclotomic kernel in
$G_{\Q,\{p\}}$ maps onto $H$, and the outer action factors
through $\Gal(\Omega_p/\Q)$.
A Levi splitting and successive conjugation put $\gamma$
in the torus: in weight $m>0$ the operator $u^m-1$ is
invertible over $\Q_p$. The commuting prime-to-$p$ action
then lies in the same torus, since no nonzero positive weight
centralizes $\gamma$. In these coordinates
\[
           \gamma s_m=u^ms_m,\qquad
           \delta s_m=\omega(\delta)^m s_m.
\]
Choose a nonzero annihilator of $D$ in $\Lambda$. It has
only finitely many roots $u^m$ with odd $m\geq3$.
Pass to the quotient of the completed free Lie algebra by
the closed ideal generated by the letters at those finitely
many weights. The quotient is the completed free Lie algebra
on the remaining letters, with the same torus action, and
the image of $H$ remains dense in every finite weight quotient.
Every odd residue class still contains infinitely many of
these weights. Let $\rho$ denote the resulting representation
and $K=\ker(\rho_p)\cap H$ the kernel of the original
pro-$p$ outer action. Functorial rational completion and
this quotient give $K\subseteq\ker\rho$.

For each remaining weight $m$, the coefficient of $s_m$
is an additive functional on $M$. It vanishes on $T$ by
parity and on $D$ because its annihilator specializes to
a nonzero scalar at $\gamma=u^m$. Thus it factors through $W$.
In the coordinates \eqref{eq:Wfree}, the coefficient has
the form
\[
              b_m\,F_a(u^m-1),\qquad a\equiv m\pmod{p-1}.
\]
Here $b_m\ne0$: otherwise density in the abelian weight-$m$
quotient would fail. The argument only uses these remaining weights.

For $h\in K$, every such coefficient is zero. A nonzero
bounded-denominator series in one variable has only finitely
many zeros in $p\Z_p$, by Weierstrass preparation. The
infinitely many weights in each odd residue class therefore
force every $F_a$ to vanish. Thus the image of $h$ in $W$
is zero.

